% 01-introduction.tex — OpenBench Software Design Document, short edition (condensed from sections/01-introduction.tex on 2026-10-05).
\section{Introduction}\label{sec:intro}

\subsection{Purpose and Scope}

This document specifies the design of Release~1 of \SRSproduct{}, the internal marketplace in which a Manager posts one Request for people, Exclusive resources, and Metered capacity, and the system allocates each Request Line under the policy of its Allocation Class and settles it in credits.

The design covers every Release~1 requirement of the SRS. It divides the system into nine back-end modules and a web client (Section~\ref{sec:arch}) and specifies each module down to Java signatures, owned tables, invariants, and failure behavior (Section~\ref{sec:modules}). The six Stretch requirements are designed only as extension points (Section~\ref{sec:evo}).

\subsection{Requirements Baseline and Changes}

The baseline is the \SRSproduct{} Software Requirements Specification, version~\SRSversion{} of 2026-10-05. Identifiers are those of its short edition: 41 Release~1 and 6 Stretch functional requirements in dot notation, twelve use cases, and the CO, UI, SI, CI, PE, SE, QA, UD, AS, and DE identifiers. Version~\SRSversion{} differs from version~0.2, the edition that was graded, by the three requirements changes below.

{\small\renewcommand{\arraystretch}{1.1}
\begin{tabularx}{\textwidth}{@{}L{2.3cm}YL{4.6cm}@{}}
\toprule
\textbf{Change} & \textbf{What changed, and why} & \textbf{Where it lands} \\
\midrule
\rc{1} One Loan at a time &
A Member with a Loan overlapping any day of a Matched line's window is not a Candidate for that line; the Shortlist is ordered by fewest person-days lent in the current Period, ties by name (\req{Profile.Availability}, \req{Matched.Shortlist.Build}). Grader feedback. &
\code{MatchedAllocator} (Section~\ref{mod:allocation}); the \code{loan} constraint (Section~\ref{sec:data}). \\
\rc{2} Inspection participation &
The inspection record of the SRS states that all four team members read and inspected it. Grader feedback. &
No design effect. \\
\rc{3} Test database &
Automated tests run on PostgreSQL~16 in a container started by the test suite (Testcontainers), with the same migrations as production; the in-memory H2 database is dropped (\req{CO-db}, \req{SI-db}). Design decision (\adr{0006}). &
Sections~\ref{sec:data} and~\ref{sec:cross}. \\
\bottomrule
\end{tabularx}
}

\subsection{Design Principles}

\textbf{Simplicity first.} The team's standing rule is the simplest architecture and program that meets the SRS: one process, one database, direct method calls, no message broker. \textbf{Modularity, low coupling, and high cohesion.} Each of the nine modules hides its entities and tables behind its application services; dependencies point down a fixed layer order with no cycles, and each module owns one cluster of the vocabulary: \module{Allocation} decides all capacity and \module{Ledger} all credit arithmetic, and neither calls the other. \textbf{Separation of concerns.} What state a Request is in (\module{Requests}), whether capacity exists (\module{Allocation}), and what it costs (\module{Ledger}) are answered by different modules. \textbf{Correctness enforced in the database.} Invariants that must survive concurrent users are database constraints, row locks, or grants, not only application checks (\adr{0006}). \textbf{Uniformity of patterns.} Every module has the same four inner layers, every user or job action is one transactional service method, every module is specified with the same template, and every error is a problem response with a stable code.

\subsection{Notation and Conventions}

The figures are UML~2 diagrams drawn with PlantUML from sources kept with this document; color carries no meaning, so every figure reads in grayscale. Three elements go beyond plain boxes and arrows; messages in sequence diagrams are numbered and named after the operations of Section~\ref{sec:modules}. \textbf{Interface lollipop and socket:} in the component diagram the port \code{ReductionGuard} is drawn as a provided interface (a ball, or lollipop) on \module{Allocation} and a required interface (a socket) on \module{Inventory}; the call at run time goes from Inventory to Allocation, but the compile-time dependency goes the other way. \textbf{Dashed arrows labeled read-only:} the dashed arrows from \module{Dashboard} stand for SQL reads through views over other modules' tables, where a solid arrow is a call to a module's application service. \textbf{\code{alt} and \code{opt} frames:} in sequence diagrams an \code{alt} frame holds alternative branches, each with its guard in brackets, and failure branches are always drawn as an \code{alt} branch.

Requirement identifiers are set in bold: functional requirements keep their dot notation, for example \req{Lifecycle.States}, and other identifiers print with the numbers the SRS gives them, for example \req{QA-correct}. Signatures are Java~17; every application-service method runs in one transaction unless marked \emph{(read-only)}. All windows, Price Bands, and Periods are in the company's one time zone.
